\rulesection{9.0 Reaction}

Leaders may attempt reaction movement during the enemy Movement Phase.

\rulesubsection{9.1 Reaction Zones}

Each leader has a reaction zone extending in all directions for a number of hexes equal to his initiative rating, including the hex he occupies.

\textbf{9.11} Reaction zones are entirely separate from ZOCs. All leaders exert reaction zones; there is no MP cost to enter one.

\textbf{9.12} Reaction zones do not extend into or through ridge hexes. They extend across a river hexside only if the hexside contains an intact bridge.

\rulesubsection{9.2 Triggering Reaction}

When any enemy unit enters a reaction zone, the player may attempt reaction movement immediately, before the enemy spends more MPs.

\textbf{9.21} The player indicates the reacting stack and rolls one die. If the result is equal to or less than the initiative of the senior leader in the stack, the reaction succeeds.

\textbf{9.22} If reaction succeeds, any or all units of that stack may immediately move as a stack up to 50\% of the basic movement allowance of the slowest type of unit in the force. Compute stacking effects first and then halve the result.

\textbf{9.23} If an enemy unit is in the reacting stack's hex, units wishing to leave must use the procedure for leaving an enemy-occupied hex. If they fail, they may not leave and may not attempt reaction again during that Movement Phase while enemy units occupy the hex.

\rulesubsection{9.3 Number of Attempts}

A leader may attempt reaction and make reaction moves any number of times in the same Movement Phase, though a leader who enters an enemy-occupied hex during reaction may not move again that phase.

\textbf{9.31} Any number of enemy leaders may react after each hex of friendly movement. A player must declare all attempts before resolving any of them.

\textbf{9.32} A player may attempt reaction any number of times against the same enemy force.

\textbf{9.33} Each stack may attempt reaction only once per hex of enemy movement.

\rulesubsection{9.4 Stacks and Reaction}

A player may attempt reaction with a stack only if enemy units have entered the reaction zone of the senior leader in the stack. Use the senior leader's initiative rating to resolve the attempt. Units beginning a reaction move in the same hex must move as a stack, though some may stay behind.

\rulesubsection{9.5 Reaction from Enemy-Occupied Hexes}

Units in an enemy-occupied hex may attempt reaction. If they succeed, they must still carry out the usual procedure to leave the hex. If they fail either attempt, they may not try again in the same phase while enemy units occupy the hex.

If a unit makes a reaction move while stacked with enemy units that entered the hex in that same turn, it may not leave through any hexside by which those enemy units entered.

\rulesubsection{9.6 Other Restrictions}

\textbf{9.61} Leaders may not use reaction in their own Movement Phase and may not react to reaction movement.

\textbf{9.62} Reaction may be attempted only after an enemy force has moved, not before it begins moving.

\textbf{9.63} Entraining within a reaction zone does not allow reaction attempts.

\textbf{9.64} Reaction may not be attempted during the Combat Phase.

\textbf{9.65} Units may not entrench or destroy or repair things using reaction movement.
